\section{Field Extensions, Cont'd}

\begin{definition}
  A field extension \( K \hookrightarrow L \) is of \underline{finite type} if there exists \(
  a_1,\ldots, a_n \in L \), s.t. \( L = K(a_1,\ldots,a_n) \). 
\end{definition}

\begin{eg}
  See that \( K(x) \) is of finite type, but not algebraic.
\end{eg}

\begin{proposition}
  \label{prop:finite-extension-is-finite-type}
  If \( K \hookrightarrow L, L \hookrightarrow L' \) are finite field extensions \( \implies K
  \hookrightarrow L' \) is finite. And see that 
  \[
    [L': K] = [L:K] \cdots [L':L]
  \] 
\end{proposition}

\begin{proof}
  Let \( e_1,\ldots, e_m \) be basis of \( L \) as \( K \)-vector space, and \( e_1',\ldots,e_n' \)
  be basis of \( L' \) as \( K \)-vector space.

  \begin{claim}
    \( \{e_i e_j' \; | \; 1 \leq i \leq m, 1\leq j \leq n\} \) be basis of \( L' \) over \( K \).
  \end{claim}

  \begin{proof}
    \leavevmode 
    \begin{itemize}
      \item System of generators: Let \( u \in L' \implies \)
        \[
          u = \sum_{j=1}^n a_j e_j' \text{ for some } a_1,\ldots, a_n \in L
        \] 
        with 
        \[
          a_j = \sum_{i=1}^m c_{ij} e_i, \; c_{ij} \in K
        \] 
        thus 
        \[
          u = \sum_{i,j} c_{ij}(e_i e_j')
        \] 
      \item Linearly independent: If we have 
        \[
          \begin{aligned}
            \sum_{i,j}\alpha_{ij}e_ie_j' &= 0, \; \alpha_{ij} \in K \\ 
            \implies \sum_{j=1}^n \bace{\sum_{i=1}^m \alpha_{ij}e_i} e_j' &= 0 \\ 
            \implies \sum_{i=1}^m \alpha_{ij} e_i &= 0 \; \forall \; i \\ 
            \implies \alpha_{ij} &= 0 \; \forall \; i,j
          \end{aligned}
        \] 
    \end{itemize}
  \end{proof}
\end{proof}

\begin{proposition}
  \label{prop:finite-type-algebraic-is-finite}
  If \( K \hookrightarrow L \) be field extension, \( a_1,\ldots,a_n \in L \) be algebraic over \( K  \implies K 
  \hookrightarrow K(a_1,\ldots, a_n)\) is finite. In particular: finite type + algebraic field
  extension is finite.
\end{proposition}

\begin{proof}
  We proceed the proof by induction on \( n \geq 1 \).
  \begin{itemize}
    \item Base case: \( n = 1 \): It is ok as we saw last time that if \( f \) be minimal polynomial
      of \( a_1 \) over \( K \), then 
      \[
        K \hookrightarrow K(a_1) = K[a_1] \cong \qo{K[x]}{(f)} \implies [K(a_1): K] = \deg(f)
      \] 
    \item Inductive case: If \( n \geq 2 \), we simply consider:
      \[
        \underbrace{K \hookrightarrow K(a_1,\ldots,a_{n-1})}_{\text{finite by induction}}
        \hookrightarrow K(a_1,\ldots, a_n)
      \] 
      with the latter one being finite by the case of \( n=1 \) and the fact that an algebraic
      element over \( K \) is clearly algebraic over the larger field \( K(a_1,\ldots,a_{n-1}) \).
      Then \textbf{Proposition} \ref{prop:finite-extension-is-finite-type} \( \implies
      \) QED.
  \end{itemize}
\end{proof}

\begin{proposition}
  \label{prop:algebraic-extension-is-algebraic}
  If \( K \hookrightarrow L \) and \( L \hookrightarrow L' \) be algebraic field extension \( \implies K \hookrightarrow 
  L'\) is an algebraic field extension.
\end{proposition}

\begin{proof}
  Suppose \( a\in L' \implies \exists \; f\in L[x] \), with \( f = x^n + b_1 x^{n-1} + \cdots + b_n \), s.t. 
  \( f(a) = 0 \). Consider 
  \[
    K \hookrightarrow K(b_1,\ldots, b_n) \hookrightarrow K(b_1,\ldots, b_n, a)
  \]
  with the first extension being algebraic by \textbf{Proposition} above
  \ref{prop:finite-type-algebraic-is-finite} and the second extension being
  algebraic since \( f \in K(b_1,\ldots, b_n)[x] \). Now \textbf{Proposition}
  \ref{prop:finite-extension-is-finite-type} \( \implies K \hookrightarrow
  K(b_1,\ldots, b_n,a) \) finite \( \implies a \) is algebraic over \( K \).
\end{proof}

\begin{proposition}
  For every field extension \( K \hookrightarrow L \), if
  \[
    L' = \{a\in L \; | \; a \text{ is algebraic over } K \}
  \] 
  \( \implies L' \) is a subfield of \( L \) containing \( K \) and \( K \hookrightarrow L' \) algebraic.
\end{proposition}

\begin{proof}
  See that \( K \subseteq L' \) is clear. In particular, \( 1 \in L' \). We need to show if \(
  a,b\in L' \), then 
  \[
    \begin{cases}
      a-b \in L' \\ 
      ab \in L' 
    \end{cases}
    \implies L' \text{ is a subring}.
  \] 
  Now apply \textbf{Proposition} \ref{prop:finite-type-algebraic-is-finite} \(
  \implies K \hookrightarrow K(a,b) \) be finite extension. And in particular, it is algebraic \(
  \implies ab, a-b \in L' \). Also, if \( a\in L' \) and \( a\ne 0 \), then \( a^{-1}\in K(a) \)
  since \( a \) is algebraic then \( K(a) \) is finite over \( K \), \( \implies a^{-1} \in L' \).
  Hence \( L' \subseteq L \) be subfield.
\end{proof}

\begin{eg}
  \leavevmode 
  \begin{enumerate}
    \item Let \( d\in \ZZ \) not a square, have 
      \[
        \QQ \hookrightarrow \QQ(\sqrt{d})
      \] 
      See that \( (\sqrt{d})^2 = d \implies \) if \( f(x) = x^2 - d \implies f(\sqrt{d}) = 0
      \implies \sqrt{d} \) is algebraic over \( \QQ \) with \( \sqrt{d} \not\in \QQ \implies f \) is
      minimal polynomial of \( \sqrt{d} \implies \QQ(\sqrt{d}) = \QQ[\sqrt{d}] \) is a degree 2
      extension of \( \QQ \).
    \item Let \( \QQ \hookrightarrow \CC \). Now we have 
      \[
        \overline{\QQ} = \{a\in \CC \; | \; a \text{ algebraic over } \QQ\} \subseteq \CC
      \] 
      See that \( \overline{\QQ} \) is countable since \( \{f \in \QQ[x]\} \) is countable and each
      polynomial has only finitely many roots in \( \CC \). Since \( \CC \) is uncountable, then 
      ``most'' element of \( \CC \) are not in \( \overline{\QQ} \).

      \begin{note}
        \( \overline{\QQ} \) is again algebraically closed, now if \( f \in \overline{Q}[x] \implies f  \)
        has a root \( a\in \CC \) since \( \CC \) is algebraically closed. Then 
        \[
          \QQ \underbrace{\hookrightarrow}_{\text{algebraic}} \overline{\QQ} \underbrace{\hookrightarrow}_{\text{algebraic}} 
          \overline{\QQ}(a) \xRightarrow{\text{prop \ref{prop:algebraic-extension-is-algebraic}}}
          a \text{is algebraic over } \QQ \implies a\in \overline{\QQ}
        \] 
      \end{note}
    \item Cyclotomic extensions of \( \QQ \): if \( n \in \ZZ_{>0} \) and \( \omega \in \CC \) be
      primitive root of \( 1 \), for example, \( \cos(\frac{2\pi}{n}) + i \sin (\frac{2\pi}{n}) \).
      \[
        \QQ \hookrightarrow \QQ(\omega)
      \] 
      is algerbaic and hence finite since \( \omega^n - 1 = 0 \) and of course not the minimal polynomial.
  \end{enumerate}
\end{eg}

\begin{definition}[\textbf{Morphisms of Field Extensions}]
  A morphism of \( K \)-extensions is a morphism of \( K \)-algebra: a ring homomorphism \( \vp :
  L_1 \to L_2 \), s.t. \( \vp \circ i_1 = i_2 \).
  \[
 		\begin{tikzcd}
			K \arrow[r, "i_1", hook] \arrow[d, "i_2"', hook] & L_1 \arrow[ld, "\vp", dashed] \\
			L_2 &
		\end{tikzcd}
  \] 
  \begin{note}
    Such \( \vp \) is always injective, hence it is an isomorphism iff it is surjective.
  \end{note}
\end{definition}

\begin{eg}
  Suppose we have 
  \[
    \begin{aligned}
      K &\hookrightarrow L_1 \\ 
      K &\hookrightarrow L_2
    \end{aligned}
  \] 
  with \( f\in K[x] \) irreducible, and \( a_1\in L_1, a_2 \in L_2 \), s.t. \( f(a_1) = 0 = f(a_2) \), then 
  \( K(a_1) \cong K(a_2) \) as extensions of \( K \).

  \begin{proof}
    Since \( f \) is irreducible, then \( f \) is the minimal polynomial of \( a_1,a_2 \) over \( K \), we saw:
    \[
      K(a_1) \cong \frac{K[x]}{(f)} \cong K(a_2)
    \] 
    which completes the proof.
  \end{proof}
\end{eg}

\section{The Galois Group}

\begin{definition}[\textbf{Galois Group}]
  If \( \qo{L}{K} \) is a field extension, the \underline{Galois Group} of \( L \) over \( K \) is 
  \[
    G\bace{\qo{L}{K}} = \{\sigma: L \to L \; | \; \sigma \text{ automorphism of field extension}\}
  \] 
  i.e. \( \sigma \) automorphism and \( \sigma(u) = u \; \forall \; u \in K \), with the group
  operation being the usual function composition.
\end{definition}

\begin{remark}
  If \( \qo{L}{K} \) finite, then any \( \sigma: L \to L \) being ring homomorphism, s.t. \(
  \sigma|_{K} = \Id \) is automatically an automorphism.
\end{remark}

\begin{proof}
  See that \( \sigma (\lambda u ) = \sigma(\lambda) \cdot \sigma(u) = \lambda \sigma(u) \) with \(
  \lambda \in K \implies \sigma \) linear over \( K \). Since \( \sigma \) is injective
  autromatically and \( \dim_{K}(L) < \infty \implies \sigma \) is also surjective.
\end{proof}

\begin{eg}
  Let \( \RR \hookrightarrow \CC \), we look into the Galois group \( G(\qo{\CC}{\RR}) \): if \(
  \sigma \in G(\qo{\CC}{\RR}) \), see that 
  \[
    \begin{aligned}
      i^2 + 1 = 0\implies \sigma(i)^2 + 1 &= 0 \\ 
      \implies \sigma(i) &= \pm i \\ 
      \implies \sigma(i) = i \implies \sigma &= \Id \\ 
      \sigma(i) = -i \implies \sigma &= \text{conjugation}
    \end{aligned}
  \] 
  Thus
  \[
    G\bace{\qo{\CC}{\RR}} = \{\Id, \text{conjugation}\}
  \] 
\end{eg}

